\documentclass[UTF8,a4paper,11pt,openany]{ctexbook}
\usepackage{../common/statlearnbook}
\SLSetBookMetadata{第 5 册}{统计学习方法讲义 v2.6.0}{采样方法、主题模型与图排序}
\begin{document}
\setcounter{page}{740}
\setcounter{chapter}{31}
\setcounter{figure}{8}
\pagestyle{headings}
\chaptermark{蒙特卡罗方法与直接采样}
\sectionmark{估计误差与诊断}
% v2.3.1 figure UID: FIG-P583-01
% Proposed post-v2.3.1 polish for FIG-P583-01: protect key-label size and stroke hierarchy.
\tikzset{slfig-FIG-P583-01/.style={font=\fontsize{9.2pt}{11pt}\selectfont,
  every path/.append style={line cap=round,line join=round}}}
\begin{figure}[htbp]
\centering
\begin{tikzpicture}[slfig-FIG-P583-01]
\begin{loglogaxis}[width=10.2cm,height=5.7cm,xmin=1,xmax=1024,ymin=.025,ymax=1.2,
  axis lines=left,xlabel={样本量 $N$},ylabel={RMSE},
  xtick={1,4,16,64,256,1024},ytick={.03125,.0625,.125,.25,.5,1},
  xticklabels={$1$,$4$,$16$,$64$,$256$,$1024$},
  yticklabels={$1/32$,$1/16$,$1/8$,$1/4$,$1/2$,$1$},
  axis line style={line width=.65pt},tick style={line width=.55pt},
  tick label style={font=\fontsize{8.6pt}{10.3pt}\selectfont},
  label style={font=\fontsize{9.6pt}{11.5pt}\selectfont},clip=false]
  \addplot[draw=SLBlue,line width=.98pt,domain=1:1024,samples=301]{x^(-.5)};
  \node[anchor=south west,font=\fontsize{9.6pt}{11.5pt}\selectfont,text=SLBlue]
    at (axis cs:120,.10) {$O(N^{-1/2})$};
  \draw[draw=SLGold!92!black,line width=.62pt] (axis cs:16,.25)--(axis cs:64,.25)--(axis cs:64,.125)--cycle;
  \node[font=\fontsize{9.2pt}{11pt}\selectfont,text=SLGold!88!black,align=center,
    yshift=2pt,fill=white,fill opacity=.94,text opacity=1,
    inner xsep=1.2pt,inner ysep=.6pt]
    at (axis cs:32,.34) {样本量 $\times4$\\误差约 $\div2$};
  \node[draw=SLRuleGray!85!black,line width=.55pt,rounded corners=2pt,fill=white,
    align=center,font=\fontsize{9.2pt}{11pt}\selectfont,inner xsep=4pt,inner ysep=3pt]
    at (axis cs:9,.055) {理论速率条件\\iid 且方差有限};
\end{loglogaxis}
\end{tikzpicture}
\caption{独立同分布且方差有限时的理论均方根误差速率；相关样本或无限方差情形不能直接沿用这条直线}\label{fig:V5-C02-rmse-rate}
\end{figure}

\begin{example}[稀有事件下的重要性抽样高方差]
在状态空间 $\{0,1\}$ 上，目标分布为 $p=(0.99,0.01)$，提议分布为 $q=(0.9999,0.0001)$，且 $h(0)=0,h(1)=100$。求 $I=E_p h$、普通重要性贡献的方差，并解释为什么少量样本可能给出误导性诊断。
\end{example}
\textbf{提示。}经验 ESS 只能描述已经观察到的权重集中度，不能发现尚未命中的稀有高贡献区域。
\end{document}
